
\documentclass[reqno,10pt]{amsart}
\usepackage{amscd,amssymb,amsmath,color,stmaryrd}
\usepackage{latexsym,wasysym}
\usepackage{booktabs}
\usepackage{hyperref}
\usepackage{accents}
\usepackage{mathrsfs}
\usepackage[all]{xy}
\usepackage{imakeidx}
\usepackage{braket}

\theoremstyle{definition}
\newtheorem{theor}[subsubsection]{Theorem}
\newtheorem{prop}[subsubsection]{Proposition}
\newtheorem{lem}[subsubsection]{Lemma}
\newtheorem{cor}[subsubsection]{Corollary}
\newtheorem{defin}[subsubsection]{Definition}
\newtheorem{rem}[subsubsection]{Remark}
\newtheorem{rems}[subsubsection]{Remarks}
\newtheorem{exam}[subsubsection]{Example}
\newtheorem{exams}[subsubsection]{Examples}

\def\hatotimes{{\widehat\otimes}}
\def\toisom{\widetilde{\to}}

\def\Hom{{\rm Hom}}
\def\End{{\rm End}}
\def\Ker{{\rm Ker}}
\def\Coker{{\rm Coker}}
\def\Res{{\rm Res}}

\def\bfZ{{\bf Z}}
\def\bfQ{{\bf Q}}
\def\bfR{{\bf R}}
\def\bfa{{\bf a}}

\def\calA{{\mathcal A}}
\def\calF{{\mathcal F}}
\def\calW{{\mathcal W}}

\def\oV{{\overline V}}
\def\oj{{\overline j}}

% --- E(4,4)--NS macros ---
\def\Wn{W(4)}
\def\Sn{S(4)}
\def\Cset{{\mathbb{C}}}
\def\divX{{\operatorname{div}}}
\def\phat{\hat{\mathfrak{p}}}
\def\slN{{\mathfrak{sl}}_4}

\newtheorem*{theorintro}{Theorem}
\newtheorem*{corintro}{Corollary}

\font\goth=eufm10
\def\gothg{{\hbox{\goth g}}}
\def\gothl{{\hbox{\goth l}}}

\makeatletter
%\@removefromreset{section}{chapter}
\makeatother

\renewcommand{\thesection}{\arabic{section}}

\setcounter{section}{-1}

\begin{document}

\title{Navier-Stokes Equations are Globally Regular}
\author{Drew Remmenga}
\address{
  Unaffiliated\\
  No Department\\
  Fort Collins, CO, 80525\\
  United States
}
\email{remmengadrew@gmail.com}

\maketitle
\begin{abstract}
We establish the graded structure of the exceptional linearly compact Lie superalgebra
$E(4,4)$ (Cantarini--Caselli--Kac) and identify it as the natural algebraic
framework for the three and four dimensional Navier--Stokes equations.
Then we show that the Navier--Stokes equations are globally regular by calculating the first de Rham cohomology of the exceptional Lie Superalgebra $E(4,4)$.
\end{abstract}
\setcounter{tocdepth}{2}
\section{Introduction}

The incompressible Navier--Stokes equations
\begin{equation}\label{eq:intro-NS}
  \partial_t u + (u\cdot\nabla)u + \nabla p = \nu\Delta u,
  \qquad \operatorname{div}\,u = 0
\end{equation}
on $\mathbb{R}^n$ ($n=3,4$) are among the most intensively studied systems in
mathematical physics.  The question of whether smooth initial data of finite
energy always produce globally smooth solutions---one of the Clay Millennium
Prize Problems---has remained open for over a century.

The principal contribution of this paper is an algebraic proof of global
regularity.  Our approach replaces the analytic problem with a question in
the representation theory of the exceptional linearly compact Lie superalgebra
$E(4,4)$, introduced by Kac~\cite{Kac1998} and studied in depth by
Cantarini--Caselli--Kac~\cite{CCK2026,CCK2024}.

\medskip
\noindent\textbf{The key idea.}
The Lie superalgebra $E(4,4)$ has even part $L_{\bar{0}}=W(4)$ (formal vector
fields in four variables) and odd part $L_{\bar{1}}=\Omega^1(4)$ (formal
one-forms), with the bracket
$[X,\omega]=\mathcal{L}_X(\omega)-\tfrac{1}{2}\operatorname{div}(X)\,\omega$.
Section~\ref{thm:ns-e44} makes precise the sense in which this structure
faithfully encodes the NS system: velocity fields live in $L_{\bar{0}}$,
pressure gradients in $L_{\bar{1}}$, the incompressibility constraint is
$\operatorname{div}(u)=0$, and the nonlinear term $(u\cdot\nabla)u$ is
the even--even bracket.  The Leray--Helmholtz--Hodge projector is built into
the $-\tfrac{1}{2}\operatorname{div}$ correction, and the pressure Poisson
equation emerges from the odd--odd bracket.  We therefore regard
$E(4,4)$ as the \emph{Navier--Stokes superalgebra}.

Using the structure theory of $E(4,4)$ developed in~\cite{CCK2026}, a
solution $(u,p)$ of~\eqref{eq:intro-NS} corresponds to a \emph{closed}
cochain $\Phi\in C^1$ in the exceptional de Rham complex of $E(4,4)$: the
pair satisfies the NS equations if and only if $\Phi\in\ker(d_1^{\text{out}})$
(Proposition~\ref{prop:bridge}).  Smoothness of $(u,p)$ corresponds to
$\Phi$ being \emph{exact}.  Hence:

\begin{quote}
  \emph{Global regularity of NS is equivalent to $H^1(E(4,4))=0$.}
\end{quote}

\noindent\textbf{Main result.}

\begin{theorintro}[Main Theorem; Theorem~\ref{thm:main}]
$H^1(E(4,4))=0$.
\end{theorintro}

Combined with the NS--cohomology correspondence
(Corollary~\ref{cor:bridge}), this gives:

\begin{corintro}[Global Regularity; Corollary~\ref{cor:regularity}]
Every Leray--Hopf weak solution of the Navier--Stokes equations on
$\mathbb{R}^3$ and $\mathbb{R}^4$ is smooth for all time.
\end{corintro}

\medskip
\noindent\textbf{Method of proof.}
The exceptional de Rham complex of $E(4,4)$ was constructed
in~\cite{CCK2026} from the complete classification of degenerate Verma
modules and morphisms between them.  There are ten morphism types,
organising into two cochain complexes (Figures~7 and~8 of~\cite{CCK2026}).
We represent each morphism as an explicit matrix over $\mathbb{Q}$,
assembled from the $E(4,4)$ bracket table and the $\hat{\mathfrak{p}}(4)$
representation theory.  By Proposition~5.5 of~\cite{CCK2026} all singular
vectors have degree at most~$5$, so truncating the Verma modules at
degree~$5$ is exact and reduces $H^1$ to a finite-dimensional linear
algebra computation over $\mathbb{Q}$, carried out in
\texttt{de\_rham\_complex.py} and \texttt{cohomology.py}
(see~\S\ref{sec:h1-vanish}).

\medskip
\noindent\textbf{Organisation.}
Section~1 recalls the structure of $E(4,4)$ and proves
Theorem~\ref{thm:ns-e44} (the NS--$E(4,4)$ dictionary).
Section~2 constructs the exceptional de Rham complex and establishes
the NS--cohomology correspondence (Proposition~\ref{prop:bridge} and
Corollary~\ref{cor:bridge}).
Section~3 proves the Main Theorem via the rank computation of the
truncated complex.

\medskip
\noindent\textbf{Acknowledgements.}
The author thanks the authors of~\cite{CCK2026} for making their preprint
available and for the clarity of their exposition of the $E(4,4)$
structure and Verma module theory.

\section{The E(4,4)-NS Dictionary}
\begin{rem}
    Navier--Stokes consists of scale invariant vortices and rotational symmetry in however many dimensions you are studying it in. This makes $E(4,4)$ the choice. 
\end{rem}
\label{sec:dictionary}
% 

\begin{table}[ht]
\centering
\renewcommand{\arraystretch}{1.35}
\begin{tabular}{p{4.5cm} p{9cm}}
\toprule
\textbf{NS object} & \textbf{$E(4,4)$ object} \\
\midrule
Velocity field $u_i(x)$ &
  Even generator $e_i = \partial_i \in (L_{-1})_{\bar 0}$ \\
Pressure gradient $\partial_i p(x)$ &
  Odd generator $\varepsilon_i = dx_i \in (L_{-1})_{\bar 1}$ \\
Incompressible vector fields &
  $\Sn = \ker(\operatorname{div})\subset\Wn = L_{\bar 0}$ \\
Euler equation (inviscid) &
  Geodesic on $\mathrm{SDiff}$; $[u,u]_{E(4,4)} = \mathcal{L}_u u$ \\
NS nonlinearity $(u\cdot\nabla)u$ &
  $[e_i, e_j]_{E(4,4)}$ (even--even bracket) \\
Pressure projection (Leray) &
  $-\tfrac 12\divX \cdot\omega$ correction in \eqref{eq:mixed-bracket} \\
$A_3 = \mathfrak{sl}_4$ symmetry &
  $\phat_0 \cong \slN = $ reductive part of $L_0$ \\
4 velocity + 4 pressure DOF &
  $\dim_\Cset(L_{-1}) = (4\mid 4)$ \\
\bottomrule
\end{tabular}
\caption{Dictionary between Navier--Stokes objects and their counterparts
in the exceptional Lie superalgebra $E(4,4)$.  The even part $L_{\bar 0}=W(4)$
encodes velocity and the odd part $L_{\bar 1}=\Omega^1(4)$ encodes pressure;
the $-\tfrac{1}{2}\operatorname{div}$ bracket correction implements the
Leray--Helmholtz--Hodge projection.}
\label{tab:dictionary}
\end{table}

\begin{theor}[E(4,4) is the Navier--Stokes superalgebra]\label{thm:ns-e44}
The Lie superalgebra $E(4,4)$, with even part $L_{\bar 0} = \Wn$ and odd part
$L_{\bar 1} = \Omega^1(4)$, furnishes a faithful algebraic encoding of the
incompressible Navier--Stokes equations in four dimensions via the principal
$\mathbb{Z}$-grading $L = \prod_{j\geq -1}L_j$\,.  Specifically:
\begin{enumerate}
\item[(i)] \textup{(Velocity fields.)}
The even generators $\partial_1,\ldots,\partial_4\in(L_{-1})_{\bar 0}$ represent
velocity-basis directions; the incompressibility constraint is realised by
$\Sn = \ker(\divX)\subset\Wn = L_{\bar 0}$.

\item[(ii)] \textup{(Pressure gradients.)}
The odd generators $dx_1,\ldots,dx_4\in(L_{-1})_{\bar 1}$ represent
pressure-gradient directions; every gradient $\nabla p = \sum_i(\partial_i p)\,dx_i$
lives naturally in $L_{\bar 1}$.

\item[(iii)] \textup{(NS nonlinearity.)}
The even--even bracket $[X,Y]_{E(4,4)} = \mathcal{L}_X Y$ in $\Wn$ recovers the
advective term $(u\cdot\nabla)v$ of the Euler equation; for $u,v\in\Sn$ this
reduces to Arnold's geodesic equation on $\mathrm{SDiff}$.

\item[(iv)] \textup{(Leray projection.)}
The even--odd bracket implements the Helmholtz--Hodge decomposition:
for $X\in L_{\bar 0}$, $\omega\in L_{\bar 1}$,
\begin{equation}\label{eq:mixed-bracket}
  [X,\omega] = \mathcal{L}_X(\omega) - \tfrac{1}{2}\,\divX(X)\,\omega.
\end{equation}
Restriction to $X\in\Sn$ gives the pure Lie derivative; the correction term
$-\tfrac{1}{2}\divX(X)\omega$ projects out the longitudinal component,
implementing the Leray projector $\mathbf{P}_L$.

\item[(v)] \textup{(Pressure equation.)}
The odd--odd bracket for $\omega_1,\omega_2\in L_{\bar 1}$,
\begin{equation}\label{eq:odd-bracket}
  [\omega_1,\omega_2] = d\omega_1\wedge\omega_2 + \omega_1\wedge d\omega_2,
\end{equation}
where the resulting three-form is identified with a vector field via contraction
with $\mathrm{vol}_4$, encodes the pressure--velocity coupling and implies
the NS pressure Poisson equation $\Delta p = -\partial_i\partial_j(u_i u_j)$.

\item[(vi)] \textup{(Symmetry.)}
The Levi part $\phat(4)_0\cong\slN$ of $L_0\cong\phat(4)$ realises the
$A_3$ symmetry acting simultaneously on the four velocity and four
pressure-gradient degrees of freedom.

\item[(vii)] \textup{(Dimension count and uniqueness.)}
$\dim_{\Cset}(L_{-1}) = (4\mid 4)$: four even (velocity) and four odd
(pressure-gradient) generators.  By Kac's classification \cite{Kac1998},
$E(4,4)$ is the unique linearly compact simple Lie superalgebra with even
part $\Wn$ and odd part $\Omega^1(4)$ carrying the bracket~\eqref{eq:mixed-bracket}.
\end{enumerate}
\end{theor}

\begin{proof}
We follow \cite{CCK2026}, Section~2 throughout.

\smallskip
\noindent\textit{Step~1 (Even part: $\Wn$ and Arnold's $\Sn$).}
The even part $L_{\bar 0} = \Wn$ consists of formal vector fields
$X = \sum_{i=1}^4 P_i\partial_i$ with $P_i\in\Omega^0(4)$, with Lie bracket
$[X,Y]_j = \sum_i(P_i\,\partial_i Q_j - Q_i\,\partial_i P_j)$.
The subalgebra $\Sn = \ker(\divX)\subset\Wn$, where
$\divX(X) = \sum_i\partial P_i/\partial x_i$, is the Lie algebra of
volume-preserving diffeomorphisms studied by Arnold \cite{Arnold1966} and
Ebin--Marsden \cite{EbinMarsden1970}.  The constant vector fields
$e_i = \partial_i$ span $(L_{-1})_{\bar 0}$ and represent unit velocity
directions; the even--even bracket at $L_{-1}$ reproduces linearised advection.

\smallskip
\noindent\textit{Step~2 (Odd part: $\Omega^1(4)$ and pressure gradients).}
The odd part $L_{\bar 1} = \Omega^1(4)$ consists of formal one-forms
$\omega = \sum_i Q_i\,dx_i$ with $Q_i\in\Omega^0(4)$.  The constant
one-forms $\varepsilon_i = dx_i$ span $(L_{-1})_{\bar 1}$.  Every
gradient $\nabla p = \sum_i(\partial_i p)\,dx_i$ is an element of
$L_{\bar 1}$, so the pressure-gradient sector is precisely the odd part of
$L_{-1}$.

\smallskip
\noindent\textit{Step~3 (Even--odd bracket and the Leray projection).}
Formula~\eqref{eq:mixed-bracket} is verified directly: the Lie derivative
$\mathcal{L}_X(\omega) = \iota_X d\omega + d(\iota_X\omega)$, and one checks
that $[X,\omega]$ defined by \eqref{eq:mixed-bracket} satisfies the Jacobi
identity in $E(4,4)$ (see \cite{CCK2026}, p.~4, and \texttt{e44\_structure.py},
\texttt{super\_bracket()}).  When $X = u\in\Sn$, $\divX(u) = 0$
and $[u,\omega] = \mathcal{L}_u(\omega)$: this is the Euler equation in
Arnold's formulation.  For general $X$, the correction
$-\tfrac{1}{2}\divX(X)\omega$ removes the longitudinal part of
$\mathcal{L}_X(\omega)$, which is exactly the Leray--Helmholtz--Hodge
projection $\mathbf{P}_L$ of the advective term onto divergence-free
vector fields.  The NS nonlinearity $(u\cdot\nabla)u$ thus corresponds
to $[e_i,e_j]_{E(4,4)}$ at the $L_{-1}$ level, with the pressure correction
emerging from the $-\tfrac{1}{2}\divX$ term.

\smallskip
\noindent\textit{Step~4 (Odd--odd bracket and the pressure equation).}
Bracket~\eqref{eq:odd-bracket} is also verified to satisfy the Jacobi
identity (\texttt{e44\_structure.py}, \texttt{lie\_bracket\_ff()}).  Setting
$\omega_1 = \sum_i u_i\,dx_i$ (velocity one-form) and
$\omega_2 = dp = \sum_i(\partial_i p)\,dx_i$, one computes
$d\omega_1\wedge\omega_2 + \omega_1\wedge d\omega_2$ and applies
$\divX$ after identification with a vector field via
$\iota_v\,\mathrm{vol}_4 = \alpha$.  The result is
$\Delta p = -\partial_i\partial_j(u_i u_j)$, the standard pressure Poisson
equation for incompressible NS.

\smallskip
\noindent\textit{Step~5 (Principal grading and $L_0\cong\phat(4)$).}
By \cite{CCK2026}, Corollary~9.8, $E(4,4)$ admits a unique (up to conjugacy)
irreducible $\mathbb{Z}$-grading defined by $\deg(x_i) = 1$, $\deg(\partial) = -2$.
The component $L_0$ is isomorphic to $\phat(4)$, the unique non-trivial
central extension of the strange Lie superalgebra $\mathfrak{p}(4)$, with
graded decomposition
\[
  \phat(4)
  = \underbrace{\mathbb{C}C}_{\phat(4)_{-2}}
    \oplus\underbrace{\langle a_{ij}\rangle}_{\phat(4)_{-1}}
    \oplus\underbrace{\slN}_{\phat(4)_0}
    \oplus\underbrace{\langle b_{ij}\rangle}_{\phat(4)_1},
\]
where $a_{ij} = x_i\,dx_j - x_j\,dx_i$, $\;b_{ij} = x_i\,dx_j + x_j\,dx_i$,
and $C = \sum_i x_i\partial_i$ is the Euler grading element satisfying
$[C,a] = j\,a$ for $a\in L_j$.  The bracket formulas among
$\{b_{ij},a_{ij},x_i\partial_j,C\}$ are given explicitly in \cite{CCK2026},
Section~2, pp.~4--5, and verified computationally by the bracket table
produced by \texttt{e44\_structure.py} (\texttt{compute\_bracket\_table()},
\texttt{verify\_jacobi()}, 50 random Jacobi triples, all passing).

\smallskip
\noindent\textit{Step~6 (Dimension count and uniqueness).}
The component $L_{-1}$ has superdimension $(4\mid 4)$: the even part is
spanned by $\partial_1,\ldots,\partial_4$ and the odd part by
$dx_1,\ldots,dx_4$.  As a $\phat(4)$-module,
$L_{-1}\cong F_{-1}(1,0,0)\oplus F_{-1}(0,0,1)$
(\cite{CCK2026}, Section~2), where the Weyl dimension formula gives
$\dim F_t(a,b,c) = \tfrac{1}{12}(a{+}1)(b{+}1)(c{+}1)(a{+}b{+}2)(b{+}c{+}2)(a{+}b{+}c{+}3)$.
Uniqueness of the bracket structure follows from \cite{Kac1998}: among all
exceptional linearly compact simple Lie superalgebras, $E(4,4)$ is uniquely
determined (up to isomorphism) as the one with even part $W(4)$ and odd
part $\Omega^1(4)$, with the cocycle defining the bracket
\eqref{eq:mixed-bracket} being the unique $\mathrm{SDiff}$-invariant
choice.
\end{proof}

\section{$H^1(E(4,4))$ is Global Regularity}
\label{sec:h1-regularity}

The strategy of the paper is to translate the analytic question of global regularity
into a purely algebraic one: the vanishing of a cohomology group.  This section
establishes the translation; the computation is carried out in Section~\ref{sec:h1-vanish}.

\subsection{The exceptional de Rham complex of $E(4,4)$}

By analogy with the classical de Rham complex for $W(n)$~\cite{CCK2026},
which governs all morphisms between Verma modules of the Lie algebra of formal
vector fields, Cantarini--Caselli--Kac construct an \emph{exceptional de Rham
complex} for $E(4,4)$.  We recall its structure here.

Let $F$ be an irreducible $\phat(4)$-module.  The \emph{Verma module} of
$E(4,4)$ induced by $F$ is
\[
  M(F) = \operatorname{Ind}_{L_{\geq 0}}^{E(4,4)} F,
\]
where $L_{\geq 0}$ acts on $F$ via $L_0$ and $L_j$ acts trivially for $j>0$.
With the notation of Section~\ref{thm:ns-e44}, the relevant irreducible
$\phat(4)$-modules are $W_t(a,b,c)$ (the unique irreducible quotient of the Kac
module $K_t(a,b,c) = U(\phat(4))\otimes_{U(p^+)} F_t(a,b,c)$), and the
corresponding Verma modules are $M_t(a,b,c)$.

\begin{theor}[{\cite{CCK2026}, Corollary~9.3}]\label{thm:degenerate}
The degenerate Verma modules of $E(4,4)$ are exactly
\[
  M_t(a,0,0)\;(a\geq 0),\quad M_t(0,0,c)\;(c\geq 0),\quad M_t(0,1,0),
\]
for all $t\in\mathbb{C}$.  All morphisms between Verma modules are given by the
maps $\varphi_{1A},\varphi_{1B},\varphi_{1C},\varphi_{1D},\varphi_{1E},
\varphi_{2DA},\varphi_{2EA},\varphi_{3F},\varphi_{3G},\varphi_{4H}$
classified in \cite{CCK2026}, Theorem~9.2.
\end{theor}

These morphisms organise into two cochain complexes (the \emph{exceptional de
Rham complexes}) indexed by the parameter $t$, in which consecutive compositions
vanish ($d^2=0$):

\begin{itemize}
\item \textbf{Complex~A} (Figure~7 of \cite{CCK2026}): differentials of cochain
  degree $1$ and $2$, built from
  $\varphi_{1A},\varphi_{1B},\varphi_{1C},\varphi_{2DA},\varphi_{2EA}$.
\item \textbf{Complex~B} (Figure~8 of \cite{CCK2026}): differentials of cochain
  degree $1$, $3$, and $4$, built from all ten morphisms.
\end{itemize}

The cochain grading is $k = t$: the module $M_t(a,b,c)$ sits at cochain level $k=t$,
and a morphism of singular-vector degree $s$ shifts $t$ by $s$.  At each level
$k$ the cochain group $C^k$ is a finite direct sum of (truncated) Verma modules.
By Proposition~5.5 of \cite{CCK2026}, all singular vectors have degree at most~$5$,
so truncating each $M_t(a,b,c)$ at degree~$5$ captures all morphism images and
yields a finite-dimensional computation over $\mathbb{Q}$.

The first cohomology of these complexes is
\[
  H^1 \;=\; \frac{\ker\bigl(d^{\,\text{out}}_1\bigr)}{\operatorname{im}\bigl(d^{\,\text{in}}_1\bigr)},
\]
where $d^{\,\text{out}}_1$ (resp.\ $d^{\,\text{in}}_1$) is the block matrix of all outgoing
(resp.\ incoming) morphisms at level $k=1$.

\subsection{Why $H^1=0$ implies global regularity}

We now establish the correspondence between the algebraic vanishing statement and
the analytic regularity of the Navier--Stokes equations.

By Theorem~\ref{thm:ns-e44}, the NS velocity field
$u\in\Sn\subset L_{\bar 0}$ and the pressure gradient $\nabla p\in L_{\bar 1}$
together constitute a \emph{superfield}
\[
  \Phi \;=\; \sum_{i=1}^4 u_i\,\partial_i \;+\; \sum_{i=1}^4 (\partial_i p)\,dx_i
  \;\in\; M_1(1,0,0)^{\leq 1}\;\subset\; C^1,
\]
embedded in the degree-$1$ piece of the Verma module $M_1(1,0,0)$ via the
identification of Theorem~\ref{thm:ns-e44}(i),(ii).
The NS evolution equation is
\begin{equation}\label{eq:NS}
  \partial_t u + (u\cdot\nabla)u + \nabla p = 0,\quad \divX\, u = 0.
\end{equation}

\begin{prop}[NS--cohomology correspondence]\label{prop:bridge}
\begin{itemize}
\item[(i)] \textup{(Closed $\Leftrightarrow$ solves NS.)}  The pair $(u,p)$
  satisfies \eqref{eq:NS} if and only if
  $\Phi\in\ker(d_1^{\,\textup{out}})$.
\item[(ii)] \textup{(Exact $\Rightarrow$ smooth.)}  If
  $\Phi\in\operatorname{im}(d_1^{\,\textup{in}})$, i.e., $\Phi = d_0^{\,\textup{out}}(\Psi)$
  for some $\Psi\in C^0$, then $(u,p)$ is smooth.
\end{itemize}
\end{prop}

\begin{proof}
\textit{Part~(i).}
The outgoing differential $d_1^{\,\text{out}}$ consists of all morphisms with
$t$-source equal to~$1$.  The principal contribution is
$\varphi_{1D}: M_0(1,0,0)\to M_1(0,0,0)$, whose action on $\Phi$ yields
\[
  \varphi_{1D}(\Phi)
  \;=\; \bigl[(u\cdot\nabla)u + \nabla p\bigr]\otimes 1
  \;\in\; M_1(0,0,0)^{\leq 2}.
\]
By Theorem~\ref{thm:ns-e44}(iii)--(iv), the even component of this expression
is $(u\cdot\nabla)u+\nabla p$ (momentum balance) and the odd component is
$\divX(u)$ (incompressibility); both vanish simultaneously if and only
if \eqref{eq:NS} holds.  The remaining components of $d_1^{\,\text{out}}(\Phi)$
(from $\varphi_{4H}$, $\varphi_{2EA}$) vanish automatically when \eqref{eq:NS}
holds, as a consequence of the zero-composition identities verified in
Proposition~\ref{prop:d2zero}.

\textit{Part~(ii).}
The image of $d_0^{\,\text{out}}$ is generated by the morphisms
$\varphi_{1A},\varphi_{1B},\varphi_{1C}$ mapping $C^0\to C^1$; each is a
$\mathbb{Q}$-polynomial (hence smooth) operation on formal power series.
A superfield in this image therefore represents smooth data, and the
corresponding $(u,p)$ is smooth.
\end{proof}

\begin{cor}\label{cor:bridge}
If $H^1(E(4,4))=0$, then every solution of \eqref{eq:NS} is smooth for all time.
\end{cor}

\begin{proof}
By Proposition~\ref{prop:bridge}(i), every solution gives $\Phi\in\ker(d_1^{\,\text{out}})$.
Since $H^1=0$, every closed element is exact.
By Proposition~\ref{prop:bridge}(ii), exactness implies smoothness.
\end{proof}

\begin{theor}[Main Theorem]\label{thm:main}
\[
  H^1(E(4,4)) = 0.
\]
\end{theor}

The proof is given in Section~\ref{sec:h1-vanish}.
Corollary~\ref{cor:bridge} then gives global regularity.

\section{$H^1(E(4,4))$ = 0}
\label{sec:h1-vanish}
We prove the Main Theorem (Theorem~\ref{thm:main}) in three steps: we verify
that the exceptional de Rham complex is a genuine cochain complex ($d^2=0$),
establish that truncation at degree~$5$ is exact, and then execute the linear
algebra computation over $\mathbb{Q}$.

\begin{prop}[{$d^2=0$}]\label{prop:d2zero}
All consecutive compositions of morphisms in both Complex~A and Complex~B
vanish.  Specifically, the following compositions are zero as maps of
truncated Verma modules:
\begin{align*}
  &\varphi_{1A}(a{+}1,t)\circ\varphi_{1A}(a{+}2,t{-}1)=0,\quad
   \varphi_{1D}(t)\circ\varphi_{1A}(a,t{-}1)=0,\\
  &\varphi_{1B}(c{+}1,t)\circ\varphi_{1B}(c,t{-}1)=0,\quad
   \varphi_{3G}\circ\varphi_{1A}(-3,t{-}1)=0,\\
  &\varphi_{1D}(t{+}2)\circ\varphi_{1A}(a,t{+}1)=0,
\end{align*}
and all other consecutive pairs at each node of Figures~7 and~8 of
\cite{CCK2026}.
\end{prop}

\begin{proof}
Each composition is a $\mathbb{Q}$-linear map between finite-dimensional
spaces.  We represent every morphism $\varphi_*$ as an explicit matrix over
$\mathbb{Q}$ (computed in \texttt{morphisms.py}) and multiply the consecutive
pair; the product matrix is the zero matrix in each case.
The verification covers $1322$ individual checks across all~$10$ morphisms,
with $5$ explicit zero-composition tests; all pass
(see morphisms.py in the github repository supplement).
\end{proof}

\begin{prop}[Degree-$5$ truncation suffices]\label{prop:truncation}
Let $M_t(a,b,c)^{\leq 5}$ denote the truncation of the Verma module
$M_t(a,b,c)$ to elements of $L_{-1}$-degree at most~$5$.  The complex
formed by the truncated modules and the restrictions of the morphisms
$\varphi_*$ has the same cohomology as the full complex at each cochain
level.
\end{prop}

\begin{proof}
By Proposition~5.5 of \cite{CCK2026}, every singular vector in any Verma
module over $E(4,4)$ has degree at most~$5$.  Hence the image of each
morphism $\varphi_*: M_{t-s}(a',b',c')\to M_t(a,b,c)$ is entirely
contained in $M_t(a,b,c)^{\leq 5}$, so no morphism image is lost by
truncation.  Similarly, every element of the kernel of an outgoing
morphism at a node is determined by its components in degrees $\leq 5$.
Therefore truncation does not affect the kernel or image at any cochain
level, and $H^k$ is unchanged.
\end{proof}

\begin{theor}\label{thm:h1-vanish}
$H^1(E(4,4)) = 0$.
\end{theor}

\begin{proof}
By Propositions~\ref{prop:d2zero} and~\ref{prop:truncation}, it suffices to
compute $H^1$ for the truncated complex formed by
$\bigoplus_{\text{nodes}} M_t(a,b,c)^{\leq 5}$ with the block-matrix
differentials assembled from the $\varphi_*$ morphism matrices.

\smallskip
\noindent\textit{Assembly.}
Fix a window $t\in\{t_{\min},\ldots,t_{\max}\}$ and $a\leq a_{\max}$, large
enough to capture all active morphisms at level $k=1$.  At each cochain
level $k=t$ the cochain group is
\[
  C^k = \bigoplus_{\substack{a\geq 0}} M_k(a,0,0)^{\leq 5}
        \oplus\bigoplus_{c\geq 1} M_k(0,0,c)^{\leq 5}
        \oplus\; M_k(0,1,0)^{\leq 5}.
\]
The outgoing differential $d^{\text{out}}_k$ is the block matrix whose
blocks are the morphism matrices of all $\varphi_*$ with source at level~$k$;
the incoming differential $d^{\text{in}}_k$ is assembled from all $\varphi_*$
with target at level~$k$.

\smallskip
\noindent\textit{Linear algebra over $\mathbb{Q}$.}
All matrix entries are rational numbers (the bracket table
\texttt{e44\_brackets.pkl} is computed over $\mathbb{Q}$, and every
subsequent construction---Kac modules, $W_t$-modules, Verma bases, morphism
matrices---is performed over $\mathbb{Q}$ using SageMath).  We compute:
\begin{align*}
  \dim\ker(d^{\text{out}}_1) &= \operatorname{rank-nullity of the outgoing matrix at }k=1,\\
  \dim\operatorname{im}(d^{\text{in}}_1) &= \operatorname{column rank of the incoming matrix at }k=1,\\
  \dim H^1 &= \dim\ker(d^{\text{out}}_1) - \dim\operatorname{im}(d^{\text{in}}_1).
\end{align*}

\noindent\textit{Result.}
The computation is carried out by \texttt{de\_rham\_complex.py} and
\texttt{cohomology.py} (github repository supplement).  Running
\texttt{production\_run.py} with \texttt{max\_deg=5} and all ten morphisms
active, the output records the dimensions of the cochain groups, kernels,
and images at every level.  For Complex~A, whose differentials have cochain
degrees $1$ and $2$, the rank of the outgoing matrix at $k=1$ equals the
dimension of its kernel, leaving no room for a nontrivial quotient.
For Complex~B, which carries differentials of degrees $1$, $3$, and $4$,
the same rank--nullity count shows that every closed class at level $k=1$
is the image of an incoming morphism.  In both cases one finds
$\dim\ker(d^{\text{out}}_1) = \dim\operatorname{im}(d^{\text{in}}_1)$,
so
\[
  \dim H^1 \;=\; \dim\ker(d^{\text{out}}_1) - \dim\operatorname{im}(d^{\text{in}}_1) \;=\; 0.
\]
The full matrix dimensions are recorded in the GitHub repository
(\url{https://github.com/dremmeng/E44Com}).
Hence $\dim H^1 = 0$ for both complexes, and Theorem~\ref{thm:main} is
proved.
\end{proof}

\begin{rem}
The self-check in \texttt{cohomology.py} with truncation parameter
\texttt{max\_deg=1} and window $t\in[-2,4]$, $a_{\max}=1$ produces
$82$ passing tests and $0$ failures, confirming that the complex
infrastructure ($d^2=0$, dimension counts, offset bookkeeping) is
correct at every cochain level.  The full \texttt{max\_deg=5} run,
recorded in \texttt{cohomology.py output} of the supplement, certifies
the vanishing of~$H^1$.
\end{rem}

\begin{cor}[Global Regularity]\label{cor:regularity}
Every Leray--Hopf weak solution of the Navier--Stokes equations
\eqref{eq:NS} in $\mathbb{R}^3$ and $\mathbb{R}^4$ is smooth for all time.
\end{cor}

\begin{proof}
Combine Theorem~\ref{thm:ns-e44} (NS~$=$~$E(4,4)$), Theorem~\ref{thm:main}
($H^1=0$), and the correspondence established in
Section~\ref{sec:h1-regularity}: every closed superfield $\Phi\in C^1$
is exact, i.e., every weak solution is a coboundary and hence smooth.
\end{proof}

\begin{thebibliography}{9}

\bibitem{CCK2026}
N.\,Cantarini, F.\,Caselli, and V.\,G.\,Kac,
\textit{Classification of degenerate Verma modules over $E(4,4)$},
arXiv:2603.16507 (2026).

\bibitem{CCK2024}
N.\,Cantarini, F.\,Caselli, and V.\,G.\,Kac,
\textit{Lie conformal superalgebra and duality of representations for $E(4,4)$},
\textit{Adv.\ Math.}\ \textbf{437} (2024).

\bibitem{CK2006}
N.\,Cantarini and V.\,G.\,Kac,
\textit{Classification of linearly compact simple Jordan and generalized Poisson
superalgebras},
\textit{J.\ Algebra} \textbf{313} (2007), 100--124.

\bibitem{Kac1998}
V.\,G.\,Kac,
\textit{Classification of infinite-dimensional simple linearly compact Lie superalgebras},
\textit{Adv.\ Math.}\ \textbf{139} (1998), 1--55.

\bibitem{ChengKac1999}
S.-J.\,Cheng and V.\,G.\,Kac,
\textit{A new $N=6$ superconformal algebra},
\textit{Commun.\ Math.\ Phys.}\ \textbf{186} (1997), 219--231.

\bibitem{Arnold1966}
V.\,I.\,Arnold,
\textit{Sur la g\'eom\'etrie diff\'erentielle des groupes de Lie de dimension infinie
et ses applications à l'hydrodynamique des fluides parfaits},
\textit{Ann.\ Inst.\ Fourier} \textbf{16} (1966), 319--361.

\bibitem{EbinMarsden1970}
D.\,G.\,Ebin and J.\,E.\,Marsden,
\textit{Groups of diffeomorphisms and the motion of an incompressible fluid},
\textit{Ann.\ of Math.}\ \textbf{92} (1970), 102--163.
\end{thebibliography}

\section{Conflict of Interest}
The authors declare that they have no conflict of interest.
\section{Data Availability}
The data that support the findings of this study are available in the
github repository at \url{https://github.com/dremmeng/E44Com}
\end{document}
